\documentclass[11pt,a4paper,reqno]{amsart}
\usepackage{iftex}
\ifPDFTeX
  \usepackage[T1]{fontenc}
  \usepackage[utf8]{inputenc}
\else
  \usepackage[no-math]{fontspec}
  \setmainfont{cmunrm.otf}[
  BoldFont=cmunbx.otf,
  ItalicFont=cmunti.otf,
  BoldItalicFont=cmunbi.otf
]
\fi
\usepackage[english]{babel}
\makeatletter
\renewcommand{\@@and}{\unskip,}
\makeatother
\usepackage{amsmath,amssymb,amsthm,mathtools}
\usepackage[margin=25mm]{geometry}
\usepackage{microtype}
\usepackage{needspace}
\usepackage[unicode,colorlinks=true,allcolors=blue]{hyperref}
\hypersetup{pdftitle={Division ring extensions and a counterexample to the pure semisimplicity conjecture},
  pdfauthor={.....}}
\setlength{\emergencystretch}{2em}
\makeatletter
\def\@settitle{\begin{center}\normalfont\large\bfseries\@title\end{center}}
\makeatother
\numberwithin{equation}{section}
\newtheorem{theorem}{Theorem}[section]
\newtheorem{lemma}[theorem]{Lemma}
\newtheorem{proposition}[theorem]{Proposition}
\newtheorem{corollary}[theorem]{Corollary}
\theoremstyle{remark}
\newtheorem{remark}[theorem]{Remark}
\newcommand{\Div}{\operatorname{Div}_k}
\newcommand{\End}{\operatorname{End}}
\newcommand{\Hom}{\operatorname{Hom}}
\newcommand{\ad}{\operatorname{ad}}
\newcommand{\Q}{\mathbb Q}
\newcommand{\Z}{\mathbb Z}

\title[Division ring extensions and pure semisimplicity]
{Division ring extensions and a counterexample to the pure semisimplicity conjecture}

\author{A.~I.~Generalov}
\address{St.~Petersburg State University,
Universitetskii pr.~28,
St.~Petersburg 198504, Russia}
\email{\href{mailto:ageneralov@gmail.com}{ageneralov@gmail.com}}

\author{A.~Tursunbaev}
\address{St.~Petersburg State University,
Universitetskii pr.~28,
St.~Petersburg 198504, Russia}
\email{\href{mailto:tursunbaev.art@gmail.com}{tursunbaev.art@gmail.com}}
\pagestyle{plain}

\renewcommand{\keywordsname}{Keywords}

\makeatletter
\renewcommand{\@setkeywords}{%
  {\itshape\keywordsname:}\enspace\@keywords\@addpunct.}
\makeatother

\keywords{pure semisimplicity, division ring extensions,
hereditary artinian rings, finite representation type, Coxeter diagrams}
\subjclass[2020]{Primary 16G60; Secondary 16K40, 16D70}

\begin{document}
\begin{abstract}
We construct a family of division ring extensions and compute their left and right dimensions. In characteristic zero, we obtain a counterexample to the pure semisimplicity conjecture: a hereditary right pure semisimple ring of infinite representation type. For each prime $p\ge3$,
we construct in characteristic $p$ a hereditary artinian ring of finite representation type with Coxeter diagram $I_2(p+2)$.
\end{abstract}
\maketitle

\section{Introduction}

A ring is called right pure semisimple if every right module over it
is a direct sum of finitely generated modules.
The pure semisimplicity conjecture asserts that such a ring has only finitely many
isomorphism classes of indecomposable finitely generated right modules.
Simson~\cite{Sim} related this question to the one-sided dimensions of bimodules
over division rings. Garc\'ia~\cite[\S\,6]{Gar} and \v{S}aroch~\cite[Proposition~1.2]{Sar}
gave criteria applicable to the triangular ring
\[
 R(H,E)=\begin{pmatrix}H&E\\0&E\end{pmatrix},
\]
where $H\subset E$ are division rings and the upper right entry is the natural
$(H,E)$-bimodule $E$.

In Schofield's example~\cite[\S\,3, p.~420]{Sch} for the Coxeter
diagram $I_2(5)$, the division ring extension $F\subset E$ satisfies
\begin{equation}\label{eq:question}
 \dim({}_F E)=2,\qquad \dim(E_F)=3,\qquad
 \dim({}_E M_3(F))=\dim(M_3(F)_E)=3 .
\end{equation}
Here $M_3(F)=\End_F(E_F)$ after a right $F$-basis has been chosen,
and $E$ is embedded in this ring by left multiplication.

We sought an explicit construction of a division ring extension $F\subset E$ with the
dimensions~\eqref{eq:question} occurring in Schofield's example.
For $k=\mathbb F_3$, we obtain such a pair of two-generated division rings and express
the generators of $F$ as words in the generators of $E$ and their inverses.

Consider the following family of extensions. Let $k$ be a field, and let $A$ be the $k$-algebra
with generators $x_i$, $i\in\Z$, and defining relations (where $[a,b]=ab-ba$)
\begin{equation}\label{eq:weyl}
 [x_{i+1},x_i]=1,\qquad [x_i,x_j]=0\quad (|i-j|>1).
\end{equation}
Let $T$ denote the division ring of fractions of $A$, let $\beta$ denote the shift
$\beta(x_i)=x_{i+1}$, and put $E=T(s;\beta)$, where $sa=\beta(a)s$ for $a\in T$.
For a subset $U\subset E$, let $\Div(U)$ denote the smallest division subring of $E$
containing $k$ and $U$. Introduce the notation
\begin{equation}\label{eq:gens}
 x=x_{-1},\quad t=x_0=sxs^{-1},\quad y=x_1=s^2xs^{-2},\quad
 v_r=y^rs,\quad F_r=\Div(x,v_r)\quad(r\in\Z).
\end{equation}

\begin{theorem}\label{thm:zero}
Let $k=\Q$. The ring $R(F_1,E)$ is right pure semisimple and has infinite representation type. Moreover,
\[
 E=\operatorname{Div}_{\Q}(x,s),\qquad
 F_1=\operatorname{Div}_{\Q}(x,s^2xs^{-1}).
\]
\end{theorem}

\begin{theorem}\label{thm:prime}
Let $k$ be a field of characteristic $p\ge3$. For $1\le r<p$,
we have the decompositions
\begin{equation}\label{eq:bibases}
 E=\bigoplus_{i=0}^{r}t^i F_r
  =\bigoplus_{i=0}^{p-r} F_r t^i .
\end{equation}
If $F=F_{p-1}$, then both dimensions of $\End_F(E_F)$ over $E$ are equal to $p$.
The ring $R(F_1,E)$ has finite representation type with Coxeter diagram $I_2(p+2)$.
\end{theorem}

For $k=\mathbb F_3$, the extension $F_2\subset E$ satisfies~\eqref{eq:question}.
The construction can be described as follows.
Relations~\eqref{eq:weyl} define the initial algebra. After passing to its division ring of fractions and adjoining the shift, both division rings are generated by two elements
over $k$.
For $k=\mathbb F_5$, the ring $R(F_1,E)$ gives an example of type $I_2(7)$.

We mention the preprint~\cite{Ren}, which also constructs a counterexample to the pure semisimplicity conjecture.

\section{Powers of an element as a right basis}

The algebra on a finite interval $x_a,\ldots,x_b$ is constructed by successively
adjoining Ore indeterminates of derivation type. When
$x_{j+1}$ is adjoined, the derivation sends $x_j$ to $1$ and vanishes
on $x_a,\ldots,x_{j-1}$. It preserves the relations~\eqref{eq:weyl},
since its values on the generators are scalars. Thus,
the ordered monomials form a basis, and the algebra on this interval is a noetherian
domain. The description of the basis implies that the embeddings obtained by
enlarging the interval are injective. Any two elements of $A$ lie in one such algebra,
so both Ore conditions hold in $A$, and the division ring $T$ is defined.
The relations~\eqref{eq:weyl} are invariant under the shift.
Since $x_i=s^{i+1}xs^{-i-1}$, we have $E=\Div(x,s)$.

\begin{lemma}\label{lem:ind}
Let $r\in\Z$, $d\ge0$, and $d!(r)_d\ne0$ in $k$, where
$(r)_d=r(r-1)\cdots(r-d+1)$ and $(r)_0=1$. The elements $1,t,\ldots,t^d$
are right linearly independent over $F_r$.
\end{lemma}

\begin{proof}
Put $v=v_r$, $P_0=1$, $P_n=x_1^r\cdots x_n^r$ for $n\ge1$.
Let $B_j$ denote the division subring of $T$ generated over $k$ by
$x_i$ for $i\le j$. Conjugation by $v$ induces on $T$
an automorphism $\theta$. For
$K=\Div(\theta^m(x)\mid m\in\Z)$, we have
\begin{equation}\label{eq:union}
 K=\bigcup_{n\ge1}P_n B_{n-1}P_n^{-1}.
\end{equation}
Indeed, $\theta^m(x)=x_{m-1}$ for $m\le0$, whereas for $m\ge1$,
the equality $v^m=P_m s^m$ gives
$\theta^m(x)=P_m x_{m-1}P_m^{-1}$. For $1\le m\le n$, we may
replace $P_m$ here by $P_n$, since $x_{m+1},\ldots,x_n$
commute with $x_{m-1}$. Moreover, $P_n$ commutes with $x_i$
for $i<0$. Hence
$P_nB_{n-1}P_n^{-1}=\Div(\theta^m(x)\mid m\le n)$,
which proves~\eqref{eq:union}.

The equality $\theta(K)=K$ gives $F_r=K(v;\theta)$. Indeed,
$F_r$ contains $K$ and $v$, and distinct powers of $v$ have distinct
degrees in $s$ over $T$. Multiplying a putative relation over $F_r$ on the right
by a common denominator from $K[v;\theta]$ and comparing the coefficients of
powers of $v$, we obtain a relation over $K$. By~\eqref{eq:union},
it therefore suffices to show that
\[
 \sum_{i=0}^d t^i P_n b_i=0,\qquad b_i\in B_{n-1},
\]
implies $b_0=\cdots=b_d=0$. The derivation
$\ad_{x_{n+1}}=[x_{n+1},\,\cdot\,]$ vanishes on
$t,P_{n-1},b_i$ and sends $x_n^h$ to $h x_n^{h-1}$ for every
$h\in\Z$. Applying it zero through $d$ times and dividing by
the corresponding nonzero factors, we obtain
\begin{equation}\label{eq:system}
 \sum_{i=0}^d x_0^i P_{n-1}x_n^{r-d+j}b_i=0
 \qquad(0\le j\le d).
\end{equation}

We prove that the matrix of the system~\eqref{eq:system} is invertible. For $0\le m<n$,
put
\[
 M_m[j,i]=x_m^i x_{m+1}^r\cdots x_{n-1}^r x_n^{r-d+j},
 \qquad 0\le i,j\le d ,
\]
with the intermediate product taken to be empty when $m=n-1$.
The matrix in~\eqref{eq:system} is $M_0$.
If column $i$ has entries $w^iY_j$, replace it
by the column with entries
\[
 (\ad_w)^i(Y_j)=\sum_{h=0}^i(-1)^{i-h}\binom ih w^hY_jw^{i-h}.
\]
Simultaneously replacing all columns amounts to right multiplication by an upper
triangular matrix with unit diagonal. For $w=x_m$ and $m<n-1$,
the new entries are
\[
 (-1)^i(r)_i x_{m+1}^{r-i}
 x_{m+2}^r\cdots x_{n-1}^r x_n^{r-d+j}.
\]
Divide the columns by $(-1)^i(r)_i$, reverse their order, and multiply
each row on the left by $x_{m+1}^{d-r}$. This yields $M_{m+1}$.

Finally, for $M_{n-1}$, put $a_j=r-d+j$. The same column replacement
gives entries $(-1)^i(a_j)_i x_n^{a_j-i}$. Multiplying row $j$ on the left
by $x_n^{-a_j}$ and column $i$ on the right by $(-1)^ix_n^i$,
we obtain the scalar matrix $((a_j)_i)_{j,i=0}^d$.
Its determinant equals the Vandermonde determinant at $a_0,\ldots,a_d$,
since $(X)_i$ is a polynomial of degree $i$ with leading coefficient $1$.
The points are pairwise distinct because $d!\ne0$, which implies that $M_0$ is invertible.
\end{proof}

\begin{lemma}\label{lem:basis}
If $r\ge1$ and the characteristic of the field $k$ is zero or greater than $r$,
then $1,t,\ldots,t^r$ form a right $F_r$-basis of $E$.
\end{lemma}

\begin{proof}
The independence follows from Lemma~\ref{lem:ind} with $d=r$.
Put $U=\sum_{i=0}^r t^i F_r$. For the operator
$\Delta(b)=bx-tb$, we have $\Delta(y^hs)=h y^{h-1}s$. Hence
\[
 st^j=y^j s=\frac{j!}{r!}\Delta^{r-j}(v_r)\in U
 \qquad(0\le j\le r).
\]
The last inclusion follows from the binomial expansion of the power of
$\Delta$, since right multiplication by $x$ commutes with
left multiplication by $t$. Moreover, $xt^j=t^jx-jt^{j-1}$. Thus $xU\subseteq U$
and $sU\subseteq U$. The set $\{a\in E:aU\subseteq U\}$ is a
division subring, since left multiplication by a nonzero element of this set
is injective on the finite-dimensional right $F_r$-space $U$, and hence
bijective. This division subring contains $k,x,s$ and therefore equals $E$.
\end{proof}

\section{Stabilizers of subspaces}

\begin{lemma}\label{lem:adjacent}
If $r\ge1$ and $r+1\ne0$ in $k$, then
\begin{equation}\label{eq:adjacent}
 F_r+F_r t=F_{r+1}+tF_{r+1}.
\end{equation}
\end{lemma}
\begin{proof}
From~\eqref{eq:weyl} and~\eqref{eq:gens}, we obtain
\[
 v_rt=v_{r+1},\qquad
 tv_{r+1}=v_{r+1}x-(r+1)v_r .
\]
Put $V=F_{r+1}+tF_{r+1}$. Then $V$ is a right $F_{r+1}$-space.
Since $xt=tx-1$, left multiplication by $x$ preserves $V$.
The last two equalities imply that $v_r$ and $v_r t$ belong to $V$.
Hence $v_rV\subseteq V$. The left stabilizer of the finite-dimensional right
space $V$ is a division subring, so $F_r V\subseteq V$.
This gives the inclusion from left to right.

Now let $U=F_r+F_rt$. The equality $tx=xt+1$ implies $Ux\subseteq U$.
Furthermore, $v_{r+1}=v_rt\in U$ and
$tv_{r+1}=v_r(tx)-(r+1)v_r\in U$. Hence
$Uv_{r+1}\subseteq U$. The right stabilizer of $U$ contains $F_{r+1}$,
whence $V\subseteq U$.
\end{proof}

For $r\ge1$ and $h\ge0$, put
$L_h^{(r)}=\sum_{i=0}^h F_rt^i$.
In characteristic zero for all $h$, and in characteristic $p$ when $r+h<p$, we have
\begin{equation}\label{eq:flag}
 L_h^{(r)}=\bigoplus_{i=0}^h t^i F_{r+h}.
\end{equation}
Indeed, $L_{h+1}^{(r)}=L_h^{(r)}+L_h^{(r)}t$.
Replacing $F_{r+h}+F_{r+h}t$ by
the space $F_{r+h+1}+tF_{r+h+1}$ according to~\eqref{eq:adjacent}, we obtain the formula by induction
on $h$. The sum is direct by Lemma~\ref{lem:basis}.

Write $H=F_1$, $L_j=L_j^{(1)}$, and put
\[
 S_j=\{a\in E:L_j a\subseteq L_j\},\qquad
 D_j=\{a\in E:L_{j-1}a\subseteq L_j\}\quad(j\ge1),\qquad
 S_0=H,\quad D_0=E .
\]
\begin{proposition}\label{prop:flags}
In characteristic zero for all $j\ge0$, and in characteristic $p$
for $0\le j\le p-2$, the following equalities hold:
\[
 S_j=F_{j+1},\qquad D_j=F_{j+1}\oplus tF_{j+1}.
\]
\end{proposition}

\begin{proof}
For $j=0$, the assertion follows from Lemma~\ref{lem:basis}.
Let $j\ge1$. By~\eqref{eq:flag}, the subspace $L_j$ is closed under right multiplication by elements of $F_{j+1}$. The set $S_j$ is a division subring, since
each of its nonzero elements acts injectively, and hence bijectively, on the finite-dimensional left $H$-space
$L_j$. Thus the right index
$[S_j:F_{j+1}]_{\mathrm r}$ divides both $\dim (L_j)_{F_{j+1}}=j+1$
and $[E:F_{j+1}]_{\mathrm r}=j+2$. Hence $S_j=F_{j+1}$.

Let $0\ne a\in D_j$. Since $1\in L_{j-1}$, by~\eqref{eq:flag}
we may write $a=\sum_{i=0}^h t^i b_i$, where $b_i\in F_{j+1}$,
$b_h\ne0$, and $h\le j$. If $h\ge2$, then $t^{j+1-h}\in L_{j-1}$,
but $t^{j+1-h}a\notin L_j$ by the right linear independence of $1,t,\ldots,t^{j+1}$
over $F_{j+1}$. Thus $h\le1$. Conversely, if
$a=b_0+tb_1$ and $b_0,b_1\in F_{j+1}$, then $t^i a\in L_j$ for $0\le i<j$.
Since $L_{j-1}=\sum_{i=0}^{j-1}Ht^i$, we obtain $a\in D_j$.
\end{proof}

\begin{proof}[Proof of Theorem~\ref{thm:zero}]
By Lemma~\ref{lem:basis}, we have $E=H\oplus tH$.
On the other hand, $t^{j+1}\notin L_j$ for every $j\ge0$ by
\eqref{eq:flag} and the right linear independence of $1,t,\ldots,t^{j+1}$ over $F_{j+1}$.
Hence $1,t,t^2,\ldots$ are left linearly independent over $H$.
Finally, by Proposition~\ref{prop:flags}, we have $D_j=S_j\oplus tS_j$
for $j\ge1$. These are all the conditions of \v{S}aroch's criterion~\cite[Proposition~1.2]{Sar}
for the sequence $f_j=t^j$. By this criterion, the ring $R(H,E)$
is right pure semisimple and is of infinite representation type.
\end{proof}

\section{Finite one-sided dimensions}

\begin{lemma}\label{lem:left}
Let $\operatorname{char} k=p>0$ and $1\le r<p$.
Then $1,t,\ldots,t^{p-r}$ form a left $F_r$-basis of $E$.
\end{lemma}
\begin{proof}
Put $q=p-r$. By~\eqref{eq:flag},
\[
 L_{q-1}^{(r)}=\bigoplus_{i=0}^{q-1}t^i F_{p-1}.
\]
The right linear independence of $1,t,\ldots,t^{p-1}$ over $F_{p-1}$
implies $t^q\notin L_{q-1}^{(r)}$. Any nontrivial left linear relation
among $1,t,\ldots,t^q$ would express its highest power $t^h$
in terms of the lower powers. Multiplying this equality on the right by $t^{q-h}$,
we would obtain $t^q\in L_{q-1}^{(r)}$. This proves independence.

Each element $x_i^p$ is central in $T$ by~\eqref{eq:weyl}, and
$v_r x_i^p v_r^{-1}=x_{i+1}^p$. In particular, $t^p\in F_r$,
since $x^p\in F_r$. Let $U=\sum_{i=0}^q F_rt^i$.
The identity $t^i x=xt^i+i t^{i-1}$ gives $Ux\subseteq U$.
Since $v_r=st^r$, we have
\[
 s=v_r(t^p)^{-1}t^q\in U .
\]
Hence $t^i s=sx^i\in U$ for $0\le i\le q$, whence $Us\subseteq U$.
The right stabilizer of $U$ contains $x,s$ and coincides with $E$. Since $1\in U$,
we have $U=E$.
\end{proof}

\begin{lemma}\label{lem:mat}
If $\operatorname{char}k=p\ge3$ and $F=F_{p-1}$, then
\[
 \dim({}_E\End_F(E_F))=\dim(\End_F(E_F)_E)=p .
\]
\end{lemma}
\begin{proof}
Put $H=F_1$ and $L=L_{p-2}$.
By~\eqref{eq:flag}, we have
\[
 L=\sum_{i=0}^{p-2}Ht^i=\bigoplus_{i=0}^{p-2}t^i F .
\]
By Lemmas~\ref{lem:basis} and~\ref{lem:left}, the class of $t^{p-1}$ in $E/L$ forms a basis both on the left over $H$
and on the right over $F$.

Let $\lambda:E\to F$ be the map assigning
to each element of $E$ the coefficient of $t^{p-1}$ in its expansion in the right $F$-basis. Its kernel is $L$. The left action of $H$ on
the one-dimensional right $F$-space $E/L$ defines a homomorphism of division rings
$\nu:H\to F$, and
\[
 \lambda(hz)=\nu(h)\lambda(z)\qquad(h\in H,\ z\in E).
\]
Since $H t^{p-1}+L=E$, we have $\nu(H)=F$.
Consider $V=\Hom_F(E_F,F_F)$ with the right action of $E$
given by $(\varphi a)(z)=\varphi(az)$.
The map $a\mapsto\lambda a$ from $E$ to $V$ is injective and
semilinear on the left with respect to the isomorphism $\nu$.
Both left spaces have dimension $p$ over the respective division rings.
Consequently, $V=\lambda E$ and $\dim V_E=1$.

Evaluation on the right $F$-basis identifies $\End_F(E_F)$ as a left
$E$-space with a sum of $p$ copies of $E$. As a right $E$-space,
this ring is identified with a sum of $p$ copies of $V$, since right multiplication
by $a\in E$ corresponds to precomposition of the $p$ coordinate functionals with left
multiplication by $a$. Thus both dimensions are equal to $p$.

\end{proof}

\needspace{7\baselineskip}
\begin{proof}[Proof of Theorem~\ref{thm:prime}]
The decompositions~\eqref{eq:bibases} and the assertion on the dimensions of the matrix ring
were proved in Lemmas~\ref{lem:basis}, \ref{lem:left}, and~\ref{lem:mat}.
Put $H=F_1$. Then $\dim({}_H E)=p$, and by
Proposition~\ref{prop:flags}, we have $\dim (D_j)_{S_j}=2$ for
$0\le j\le p-2$. By Lemma~1.4 of \v{S}aroch~\cite{Sar}, the right
dimension sequence of the bimodule $ {}_H E_E$ is periodic with period
\[
 (1,\underbrace{2,\ldots,2}_{p-1},1,p).
\]
Here the terms of the sequence are the right dimensions of the
bimodule itself and its successive right dual bimodules.
By~\cite[Lemma~1.3]{Sar}, the ring $R(H,E)$ has exactly $p+2$
isomorphism classes of indecomposable finitely generated right modules.
By the definition of Dowbor, Ringel, and Simson~\cite[\S\,3]{DRS}, its
Coxeter diagram is \mbox{$I_2(p+2)$}.
\end{proof}

\begin{remark}\label{rem:three}
In the case $k=\mathbb F_3$, the one-sided relations for $t$ can
be written explicitly. Let $F=F_2$, $v=y^2s$, and $b=v x v^{-1}\in F$.
Then $b=y^2t y^{-2}=t-y^{-1}$. From $[t,y^{-1}]=y^{-2}$ it follows that
\[
 t^2+b t+b^2=0,\qquad t^3=v x^3 v^{-1}\in F .
\]
The first equality expresses $t^2$ in the left basis $1,t$.
On the right over the same division ring, $1,t,t^2$ are independent by Lemma~\ref{lem:basis}.
Thus, in~\eqref{eq:question}, one can take
$E=\Div(x,s)$, $F=\Div(x,(s^2xs^{-2})^2s)$ over the prime field
$\mathbb F_3$.
\end{remark}

\begin{thebibliography}{9}

\bibitem{Sim}
D.~Simson,
\emph{A class of potential counterexamples to the pure semisimplicity conjecture},
in: Advances in Algebra and Model Theory, Algebra, Logic and Applications,
vol.~9, Gordon and Breach, Amsterdam, 1997, 345--373.

\bibitem{Gar}
J.~L.~Garc\'ia,
\emph{Small potential counterexamples to the pure semisimplicity conjecture},
Journal of Algebra and Its Applications \textbf{17} (2018), no.~10,
\href{https://doi.org/10.1142/S0219498818501839}{1850183}.

\bibitem{Sar}
J.~\v{S}aroch,
\emph{Pure semisimplicity conjecture and Artin problem for dimension sequences},
Journal of Pure and Applied Algebra \textbf{225} (2021), no.~11,
\href{https://doi.org/10.1016/j.jpaa.2021.106745}{106745}.

\bibitem{Sch}
A.~Schofield,
\emph{Hereditary artinian rings of finite representation type and extensions
of simple artinian rings},
Mathematical Proceedings of the Cambridge Philosophical Society
\textbf{102} (1987), no.~3, \href{https://doi.org/10.1017/S0305004100067463}{411--420}.

\bibitem{DRS}
P.~Dowbor, C.~M.~Ringel, D.~Simson,
\emph{Hereditary artinian rings of finite representation type},
in: Representation Theory~II, Lecture Notes in Mathematics, vol.~832,
Springer-Verlag, Berlin, 1980, 232--241.

\bibitem{Ren}
J.~Ren,
\emph{A counterexample to the pure semisimplicity conjecture},
\href{https://arxiv.org/abs/2610.39883}{arXiv:2610.39883} (2026).

\end{thebibliography}
\end{document}